\section{Fourier Transform -- Practice}

Due to the result of the convolution theorem:
$\mathcal{F}\{f * g\} = \mathcal{F}\{f\} \cdot \mathcal{F}\{g\}$
one can apply a Gaussian filter by multiplying the Fourier transform of
the image with the Fourier transform of the Gaussian kernel.
I.e. the product in the frequency domain corresponds to convolution in the spatial domain.
Rather than taking the Fourier transform of the Gaussian kernel,
we can derive it analytically and code it directly.
Let $g(x, y) = \frac{a}{\pi} e^{-a(x^2 + y^2)}$ be a 2D Gaussian kernel with parameter $a > 0$.
The derivation is as follows:

\begin{align}
    \mathcal{F}\{g\}(u, v)
    &= G(u, v)
    = \int_{-\infty}^{\infty} \int_{-\infty}^{\infty} g(x, y) e^{-i 2 \pi (ux + vy)} \, dx \, dy \\
    &= \frac{a}{\pi} \int_{-\infty}^{\infty} \int_{-\infty}^{\infty}
       e^{-a(x^2 + y^2)} e^{-i 2 \pi (ux + vy)} \, dx \, dy \\
    &= \frac{a}{\pi} \int_{-\infty}^{\infty} e^{-ax^2 -i 2 \pi ux} \, dx
       \int_{-\infty}^{\infty} e^{-ay^2 -i 2 \pi vy} \, dy \\
    &= \frac{a}{\pi}
       \sqrt{\frac{\pi}{a}} e^{\left(
		   \frac{-\pi^2 u^2}{a}
	   \right)}
       \cdot
       \sqrt{\frac{\pi}{a}} e^{\left(
		   \frac{-\pi^2 v^2}{a}
   	   \right)} \\
    &= \frac{a}{\pi} \cdot \frac{\pi}{a}
       e^{\left(
		   \frac{-\pi^2 (u^2 + v^2)}{a}
   	   \right)} \\
    &= e^{\left(
		\frac{-\pi^2 (u^2 + v^2)}{a}
	\right)}
\end{align}
where the separability of the exponent in line 3 allows the double integral to factor
into a product of two independent 1D integrals i.e. they only depend on $x$ and $y$ respectively,
and thus can be evaluated separately.
Each integral is evaluated by completing the square in the exponent,
yielding the standard Gaussian integral result
$\int_{-\infty}^{\infty} e^{-(ax^2 + bx)} \, dx 
= \sqrt{\frac{\pi}{a}} e^{\left(
	\frac{b^2}{4a}
\right)}$
with $b = 2\pi i u$ and $b = 2\pi i v$ respectively.
\footnote{
	Note that since we square $b$ in the exponent,
	the imaginary part cancels out.
}
Substituting $a = \frac{1}{2\sigma^2}$ yields the filter used in the implementation:
\begin{equation}
    G(u, v) = \exp\!\left(-2\pi^2 \sigma^2 (u^2 + v^2)\right)
\end{equation}

\begin{lstlisting}
def scale_fft(
    image: NDArray[Any],
    sigma: float,
) -> NDArray[np.float64]:
    assert sigma > 0, "sigma must be positive"
    assert image.ndim == 2, "image must be a 2-D grayscale array"

    rows, cols = image.shape

    # Frequency coordinates in cycles-per-pixel, range [-0.5, 0.5)
    u = np.fft.fftfreq(cols)
    v = np.fft.fftfreq(rows)
    # Prepare and perform fft on image
    U, V = np.meshgrid(u, v)
    F = np.fft.fft2(image)
    # Multiply in frequency space with Gaussian filter kernel function G
    G = np.exp(-2 * np.pi**2 * sigma**2 * (U**2 + V**2))
    F_filtered = F * G

    # Inverse FFT; take real part to discard numerical imaginary residuals
    return np.real(np.fft.ifft2(F_filtered))

\end{lstlisting}

The implementation constructs the frequency coordinates using \texttt{np.fft.fftfreq},
which generates the appropriate frequency bins for the given image dimensions,
such that the Nyquist frequency corresponds to 0.5 cycles per pixel.
Based on the frequency coordinates, we can compute the Gaussian filter kernel $G(u, v)$ directly in the frequency domain,
and then apply it to the Fourier transform of the image.
Finally, the inverse Fourier transform is computed to obtain the filtered image in the spatial domain,
and the real part is taken to eliminate any numerical artifacts that may arise from the inverse transform.

\begin{figure}[H]
	\centering
	\includegraphics[width=0.8\textwidth]{images/scale_fft.png}
	\caption{Gaussian filtering in the frequency domain, with sigmas = $\{1, 3, 7, 15, 30\}$}
	\label{fig:scale_fft}
\end{figure}

The above figure \ref{fig:scale_fft} illustrates the expected blurring effect of applying a Gaussian filter.
Increasing $\sigma$ widens the Gaussian in the spatial domain, meaning each output pixel 
is a weighted average with increasingly even contribution of neighboring pixels, 
thus resulting stronger blurring.
Equivalently, in the frequency domain, a larger $\sigma$ narrows the Gaussian filter $G(u, v)$, 
attenuating high frequencies more aggressively -- which correspond to fine details and edges -- 
while preserving low frequencies, which correspond to coarse structure.
